\documentclass[12pt,a4paper]{article}
\usepackage[utf8]{inputenc}
\usepackage[T1]{fontenc}
\usepackage[brazil]{babel}
\usepackage{geometry}
\usepackage{longtable}
\usepackage{array}
\usepackage{booktabs}
\usepackage{hyperref}
\usepackage{fancyhdr}
\geometry{left=2.3cm,right=2.3cm,top=2.45cm,bottom=2.4cm}
\hypersetup{hidelinks}
\renewcommand{\arraystretch}{1.25}
\pagestyle{fancy}
\fancyhf{}
\fancyhead[L]{\small RIFTLAB / APS}
\fancyhead[R]{\small REQUISITOS FUNCIONAIS}
\fancyfoot[C]{\small\thepage}
\renewcommand{\headrulewidth}{0.4pt}
\setlength{\headheight}{15pt}
\addtolength{\parskip}{4pt}

\begin{document}
\begin{titlepage}
  \setlength{\parindent}{0pt}\raggedright
  {\normalsize\bfseries INSTITUTO FEDERAL DE EDUCAÇÃO, CIÊNCIA E\\ TECNOLOGIA DO CEARÁ\par}
  \vspace*{2.0cm}
  {\Large\bfseries ANÁLISE E PROJETO DE SISTEMAS\par}
  \vspace{1.7cm}
  {\Huge\bfseries RiftLab\par}
  \vspace{0.4cm}
  {\rule{9cm}{1pt}\par}
  \vspace{0.7cm}
  {\LARGE Especificação de requisitos funcionais\par}
  \vspace{0.5cm}
  {\large Visão final planejada do produto\par}
  \vspace{1.6cm}
  {\large Aluno: Juliano Magalhães Jurity\par}
  \vspace{0.3cm}
  {\large Professor: Thiago Queiroz de Oliveira\par}
  \vfill
  {\large RF01--RF22 \quad | \quad UC01--UC09\par}
  \vspace{0.3cm}
  {28 de setembro de 2026\par}
\end{titlepage}

\section*{Resumo da especificação}
\begin{center}
\begin{tabular}{@{}p{3.4cm}p{10.3cm}@{}}
\toprule
\bfseries Campo & \bfseries Definição \\
\midrule
Produto & Ferramenta local de análise e estudo de decisões de compra de itens após partidas concluídas de League of Legends. \\
Público & Jogadores iniciantes e de elos mais baixos; operador local e revisor. \\
Interface & Quatro áreas: Painel, Campeões, Minha posição e Histórico. A jornada de estudo começa no Histórico. \\
Fonte de verdade & Contexto anterior à compra, catálogo do patch, qualidade observada e versões persistidas. \\
Limite & Sem recomendação durante partida em andamento; sem ranking entre jogadores; demonstração sintética identificada. \\
Natureza & Requisitos da entrega final prevista, não declaração de implementação atual. \\
\bottomrule
\end{tabular}
\end{center}
\clearpage

\section{Objetivo e escopo}
O RiftLab é uma aplicação local de estudo de decisões de compra de itens em partidas
\emph{concluídas} de League of Legends. A pessoa acompanha o próprio desempenho no Painel,
abre uma partida pelo Histórico, consulta o contexto anterior à compra, considera
alternativas e, no modo Estudo, escolhe antes de ver a comparação. A área Campeões reúne
guias de estudo por campeão, e Minha posição mostra o elo oficial da própria conta. O
sistema também prevê importação, avaliação e experimentos reproduzíveis. Este documento
descreve o \emph{escopo pretendido do produto}; um requisito listado aqui não significa que
sua implementação já esteja concluída.

O produto não oferece aconselhamento durante partidas em andamento, não compara o usuário
com outros jogadores e não promete subida de elo. O modo de demonstração usa dados
sintéticos e deve ser identificado como tal.

\section{Atores}
\begin{description}
  \item[Jogador:] usa o Painel, o Histórico, os guias de Campeões e Minha posição; estuda
  decisões, define objetivo e realiza cenário de transferência.
  \item[Operador local:] configura a aplicação e a conta própria, importa partidas,
  acompanha processamento, executa experimentos e remove dados locais.
  \item[Revisor:] julga comparações de forma cega, pode se abster e exporta avaliação.
  \item[Serviço de dados da Riot:] fornece partidas concluídas (Match-V5), o elo oficial
  somente da conta própria (League-V4) e ativos estáticos versionados (Data Dragon). Não
  participa do modo de demonstração.
\end{description}

\section{Interface final}
A interface segue o protótipo conceitual do RiftLab (Painel desktop e mobile), com
navegação fixa entre quatro áreas.

\begin{center}
\begin{tabular}{@{}p{2.9cm}p{8.4cm}p{2.4cm}@{}}
\toprule
\bfseries Área & \bfseries Conteúdo & \bfseries Requisitos \\
\midrule
Painel & Taxa de vitória e KDA médio do recorte, decisões revistas, campeões no recorte e em foco, cinco partidas mais recentes e card de elo. & RF18--RF20, RF22 \\
Histórico & Lista completa de partidas com busca, filtros e resumo; cada partida abre o contexto e a jornada de estudo. & RF02--RF08, RF13--RF16, RF20 \\
Campeões & Guia de estudo por campeão com retrato oficial, habilidades, trajeto de itens observado e perguntas de matchup. & RF21 \\
Minha posição & Elo, divisão, PDL e próximo degrau da conta própria. & RF22 \\
\bottomrule
\end{tabular}
\end{center}

\section{Requisitos funcionais}
Os critérios abaixo descrevem resultados observáveis. Casos sem dados suficientes devem
exibir a limitação, sem produzir uma recomendação como se fosse confiável. Prioridade:
\textbf{E} = essencial; \textbf{I} = importante; \textbf{D} = desejável. Em RF09, o formulário é essencial e o texto interpretado é desejável.

\small
\begin{longtable}{@{}p{1.1cm}p{5.8cm}p{6.4cm}p{1.1cm}@{}}
\caption{Requisitos funcionais do RiftLab}\label{tab:rf}\\
\toprule
\bfseries ID & \bfseries Requisito & \bfseries Critério de verificação & \bfseries Prior. \\
\midrule
\endfirsthead
\toprule
\bfseries ID & \bfseries Requisito & \bfseries Critério de verificação & \bfseries Prior. \\
\midrule
\endhead
\bottomrule
\endfoot
\multicolumn{4}{@{}l}{\bfseries Importação e revisão de compras} \\
RF01 & O sistema deve permitir abrir uma demonstração sem chave de API nem conexão externa. & A jornada demonstrativa abre com dados sintéticos identificados e não faz coleta de partidas reais. & E \\
RF02 & O operador deve poder solicitar a importação de uma partida concluída por identificador BR1 válido, sem duplicar a mesma partida ou um trabalho ativo. & Identificador inválido é recusado; requisições repetidas reaproveitam o trabalho ou a partida correspondente. & E \\
RF03 & O sistema deve apresentar andamento, sucesso, falha e opção de tentar novamente quando aplicável. & A interface distingue os estados e informa falhas recuperáveis sem mostrar credenciais ou detalhes internos. & E \\
RF04 & O jogador deve poder selecionar um participante e um ponto de compra dentro do recorte suportado. & A lista associa cada decisão a participante e momento; cenários fora do recorte mostram motivo. & E \\
RF05 & O sistema deve mostrar o estado anterior à compra, a origem dos dados e as ausências ou incertezas relevantes. & Inventário, orçamento aproximado e qualidade são apresentados com indicação da fonte; dado ausente não vira valor inventado. & E \\
RF06 & O sistema deve gerar alternativas de itens permitidas pelo catálogo versionado do patch da partida. & Itens incompatíveis ou sem catálogo válido não são apresentados como candidatos legais. & E \\
RF07 & O jogador deve poder comparar até três alternativas, inclusive custo restante, sem exigir exatamente três. & A comparação funciona com uma, duas ou três alternativas e sinaliza orçamento desconhecido. & E \\
RF08 & O sistema deve ordenar alternativas com um modelo versionado e sinalizar suporte insuficiente. & A resposta identifica a versão do modelo e se abstém de ranking quando as condições de qualidade ou suporte falham. & E \\
RF09 & O jogador deve poder definir o objetivo por formulário e, opcionalmente, por texto interpretado e validado. & O formulário está sempre disponível; a interpretação do texto requer confirmação e, se ambígua ou indisponível, devolve o usuário ao formulário. & E/D \\
RF10 & O revisor deve poder registrar julgamento e exportar relatório JSON ou HTML. & A avaliação admite abstenção; o relatório inclui contexto e versões, sem credenciais ou identificadores pessoais brutos. & E \\
RF11 & O operador deve poder executar experimentos reproduzíveis e consultar o manifesto do modelo. & O experimento registra versões, dados e parâmetros necessários à repetição; o modelo ativo pode ser identificado. & E \\
RF12 & O operador deve poder expurgar uma partida local e os derivados associados. & A operação pede confirmação e informa dados, instantâneos e modelos afetados; não afirma apagar efeitos de treinamento passado. & E \\
\multicolumn{4}{@{}l}{\bfseries Modo Estudo} \\
RF13 & No modo Estudo, o sistema deve ocultar ranking e compra histórica até a ação explícita de revelar, inclusive após escolha ou resposta ``não sei''. & Antes da revelação a resposta histórica não aparece; a revelação depende de ação inequívoca. & E \\
RF14 & O modo Estudo deve oferecer pistas graduais e glossário curado. & A pessoa pode abrir as pistas separadamente; os termos apresentados têm explicação acessível. & E \\
RF15 & O sistema deve persistir tentativa, objetivo, escolha e pistas utilizadas, distinguindo comparação de gabarito absoluto. & O registro permite reconstruir a tentativa sem afirmar que existe uma única decisão correta. & E \\
RF16 & O sistema deve oferecer cenário de transferência independente, sem pistas, e exportação da reflexão. & Após a comparação, a pessoa resolve outro cenário sem ajuda e pode registrar reflexão. & I \\
RF17 & O sistema deve recusar fontes de partidas em andamento e não integrar dados do cliente ao vivo. & Solicitações de cenário ao vivo são rejeitadas; a análise usa apenas partida concluída ou demonstração sintética. & E \\
\multicolumn{4}{@{}l}{\bfseries Navegação, Painel, Histórico, Campeões e Minha posição} \\
RF18 & O sistema deve oferecer navegação entre Painel, Campeões, Minha posição e Histórico em desktop e mobile, com aviso visível quando os dados forem de demonstração. & As quatro áreas são acessíveis por menu superior (desktop) e barra inferior (mobile), por teclado, e indicam a área ativa em texto; dados sintéticos aparecem marcados. & E \\
RF19 & O Painel deve mostrar taxa de vitória e KDA médio das últimas partidas concluídas da conta própria, decisões revistas, campeões no recorte e campeões em foco. & Cada métrica mostra o tamanho da amostra (até 20 partidas, ou menos quando não houver 20); recorte vazio exibe orientação em vez de zeros. & E \\
RF20 & O Painel deve listar as cinco partidas mais recentes com acesso ao Histórico completo, e o Histórico deve permitir buscar, filtrar por resultado e abrir o contexto de cada partida. & A lista mostra resultado, campeão, data, duração, K/D/A e tropas; ``Contexto'' leva à partida e às decisões de compra (RF04). & E \\
RF21 & A área Campeões deve apresentar um guia de estudo por campeão, com nome e retrato oficiais do patch, habilidades, trajeto de itens observado e perguntas de matchup. & Nome e retrato vêm da versão fixada do patch, com texto de reserva; o trajeto mostra frequência com denominador e cobertura e se abstém quando a amostra é insuficiente; nada é apresentado como ``build ideal''. & I \\
RF22 & Minha posição deve exibir o elo oficial apenas da conta do próprio usuário: elo, divisão, PDL de 0 a 100 e próximo degrau. & A variação de PDL compara instantâneos locais e só é associada a uma partida quando não houver ambiguidade; conta não classificada ou serviço indisponível mostram estado claro; nenhuma outra conta é consultada ou exibida. & I \\
\end{longtable}
\normalsize

\section{Fluxo principal da experiência final}
\begin{enumerate}
  \item O operador abre a demonstração ou importa partidas concluídas; o sistema mostra
  progresso e permite recuperação de falhas aplicáveis (RF01--RF03, RF17).
  \item No Painel, o jogador vê o desempenho do recorte com o tamanho da amostra, as cinco
  partidas mais recentes, os campeões em foco e o card de elo (RF18--RF20, RF22).
  \item No Histórico, busca ou filtra uma partida e abre o contexto. Escolhe participante e
  compra; o sistema congela o catálogo do patch e mostra o estado anterior, sua origem e
  suas incertezas (RF04--RF06, RF20).
  \item O jogador confirma o objetivo. No modo Estudo, escolhe entre alternativas em ordem
  neutra, com pistas opcionais, sem ver compra histórica ou ranking antes da revelação
  explícita (RF09, RF13--RF14).
  \item Após a revelação, o sistema compara até três alternativas, custo restante, suporte
  e versão do modelo. Quando a evidência não permite ordenar, apresenta abstenção e motivo
  (RF07--RF08).
  \item A tentativa, a reflexão em outro cenário, a avaliação cega e as versões ficam
  disponíveis para revisão e exportação; o operador pode consultar experimentos e remover
  dados locais com confirmação (RF10--RF12, RF15--RF16).
  \item Em Campeões, o jogador estuda padrões observados de um campeão, com amostra e
  cobertura explícitas (RF21).
\end{enumerate}

\section{Regras que atravessam os requisitos}
\begin{itemize}
  \item \textbf{Tempo da decisão:} inventário, ouro e qualidade para elegibilidade são
  calculados antes do item comprado; informação posterior não pode influenciar essa
  fotografia da decisão.
  \item \textbf{Ausência e incerteza:} orçamento desconhecido, catálogo indisponível ou
  reconciliação não validada devem aparecer como limites, inclusive quando impedem
  recomendação. Valor ausente nunca deve ser apresentado como zero observado.
  \item \textbf{Aprendizagem:} ranking e compra histórica continuam ocultos após a escolha
  ou ``não sei'' até o comando explícito de revelar. A comparação não é um gabarito
  universal; o cenário de transferência deve ser independente e sem pistas.
  \item \textbf{Métricas descritivas:} taxa de vitória, KDA e frequências descrevem o recorte
  e sempre mostram o denominador; não são usadas para afirmar causa, avaliar o aluno ou
  prometer subida de elo.
  \item \textbf{Sem comparação entre jogadores:} o elo exibido é só o da conta própria; não
  existe escada, pódio ou perfil público.
  \item \textbf{Proveniência:} análise e relatório referenciam patch, catálogo, esquema,
  modelo e origem dos dados. A exclusão local não desfaz automaticamente um modelo já
  treinado; a interface deve informar essa limitação.
\end{itemize}

\section{Rastreabilidade dos casos de uso}
\begin{center}
\begin{tabular}{@{}p{1.6cm}p{5.2cm}p{3.4cm}p{3.6cm}@{}}
\toprule
Caso & Nome & Requisitos & Área da interface \\
\midrule
UC01 & Importar partida ou abrir demonstração & RF01--RF03, RF17 & Histórico \\
UC02 & Revisar compra & RF04--RF08 & Histórico \\
UC03 & Informar objetivo & RF09 & Histórico \\
UC04 & Avaliar e exportar & RF10 & Histórico \\
UC05 & Operar e remover dados & RF11--RF12 & Operação local \\
UC06 & Estudar decisão & RF13--RF16 & Histórico \\
UC07 & Consultar painel e histórico & RF18--RF20 & Painel, Histórico \\
UC08 & Estudar campeão & RF21 & Campeões \\
UC09 & Consultar minha posição & RF22 & Minha posição, Painel \\
\bottomrule
\end{tabular}
\end{center}

\section{Restrições da especificação}
A análise é pós-partida, local e de estação individual. As telas da visão final são
protótipos conceituais e usam valores ilustrativos; não constituem evidência de
integração com API externa, treinamento de modelo ou medição dos requisitos não
funcionais. Retratos (Data Dragon) e emblemas de elo (Ranked Emblems) são ativos da Riot
Games usados com crédito; o RiftLab não é endossado pela Riot Games.

\end{document}
